\newpage
\section{Counting \& Combinatorics}

\subsection{Counting by addition}
\label{sec:8-1}
% CN 8.1

\begin{theorem}[Sum rule]
If $A,B$ are disjoint finite sets, then $|A\cup B|=|A|+|B|$. More generally, for pairwise disjoint $A_1,\dots,A_m$:
\[
|A_1\cup A_2\cup\cdots\cup A_m| = |A_1|+|A_2|+\cdots+|A_m|.
\]
\end{theorem}
\begin{note}[Counting meaning]
Already established set-theoretically at \S\ref{sec:1-12}. Combinatorially: choosing one item from $m$ pairwise disjoint option-sets has $\sum_i|A_i|$ total possibilities.
\end{note}

\begin{example}
Two disjoint lists, lengths $r,c$; concatenating and summing all entries costs $r+c$ additions.
\end{example}

\bigskip
\begin{theorem}[Bijection rule (``Zeroth Rule of Counting'')]
\label{thm:bijectionrule}
If there is a bijection $f: A\to B$ between finite sets $A,B$, then
\[
|A| = |B|.
\]
\end{theorem}
\begin{note}[Counting by correspondence]
Every combinatorial argument below rests on this: to count a hard set, put it in one-to-one correspondence with an easy one. Bijections: \S\ref{sec:2-7}. Used for subsets $\leftrightarrow$ bit strings (\S\ref{sec:1-10}), derangements (\S\ref{sec:8-4}) and multisets $\leftrightarrow$ $0/1$ strings (\S\ref{sec:8-9}).
\end{note}

\subsection{Counting by multiplication}
\label{sec:8-2}
% CN 8.2

\begin{theorem}[Product rule]
For finite $A,B$: $|A\times B| = |A|\cdot|B|$. More generally, for a task done in $m$ independent steps, with $n_i$ options at step $i$:
\[
\text{total options} = n_1\cdot n_2\cdots n_m.
\]
\end{theorem}
\begin{note}[Independence of the steps]
Already established at \S\ref{sec:1-14}. ``Independent'' is essential: the count of options at step $i$ must not depend on earlier choices.
\end{note}

\begin{example}
An $r\times c$ table: summing every entry via nested loops costs $rc$ additions -- exactly $|\{1,\dots,r\}\times\{1,\dots,c\}|$.
\end{example}

\subsection{Inclusion-exclusion}
\label{sec:8-3}
% CN 8.3

\begin{proposition}[Two sets]
\[
|A\cup B| = |A|+|B|-|A\cap B|.
\]
\end{proposition}
\begin{note}[Correcting the double count]
Already proved at \S\ref{sec:1-12}. $|A|+|B|$ double-counts $A\cap B$; subtracting once corrects it.
\end{note}

\begin{proposition}[Three sets]
\[
|A\cup B\cup C| = |A|+|B|+|C| - |A\cap B|-|A\cap C|-|B\cap C| + |A\cap B\cap C|.
\]
\end{proposition}
\begin{note}[Why the correction pattern]
Summing $|A|+|B|+|C|$ counts elements in exactly two sets twice, elements in all three sets three times. Subtracting all pairwise intersections fixes the double-count but now over-corrects triple-members (removed three times, once per pair, instead of the needed two). Adding back $|A\cap B\cap C|$ restores them to count $1$.
\end{note}

\bigskip
\begin{theorem}[General inclusion-exclusion]
For finite $A_1,\dots,A_n$:
\[
\left|\bigcup_{i=1}^n A_i\right| = \sum_{k=1}^n (-1)^{k+1} \sum_{1\le i_1<\cdots<i_k\le n} |A_{i_1}\cap\cdots\cap A_{i_k}|.
\]
\end{theorem}
\begin{proof}[Proof by induction on $n$]
\textbf{Base case} ($n=1$): both sides equal $|A_1|$.

\textbf{Inductive step}: Assume the formula for $n$ sets. Write $U=A_1\cup\cdots\cup A_n$. Then
\[
\left|\bigcup_{i=1}^{n+1}A_i\right| = |U \cup A_{n+1}| = |U| + |A_{n+1}| - |U\cap A_{n+1}|
\]
by the two-set formula. Since $U\cap A_{n+1} = \bigcup_{i=1}^n (A_i\cap A_{n+1})$ (distributive law, \S\ref{sec:1-12}), the inductive hypothesis applies to \emph{both} $|U|$ and $|U\cap A_{n+1}|$ (the latter as a union of $n$ sets $A_i\cap A_{n+1}$):
\[
|U| = \sum_{k=1}^n(-1)^{k+1}\!\!\sum_{i_1<\cdots<i_k\le n}\!\! |A_{i_1}\cap\cdots\cap A_{i_k}|,
\]
\[
|U\cap A_{n+1}| = \sum_{k=1}^n(-1)^{k+1}\!\!\sum_{i_1<\cdots<i_k\le n}\!\! |A_{i_1}\cap\cdots\cap A_{i_k}\cap A_{n+1}|.
\]
The first sum ranges over all size-$k$ subsets of $\{1,\dots,n\}$ (excluding $A_{n+1}$); the second, negated, ranges over all size-$(k{+}1)$ subsets of $\{1,\dots,n+1\}$ that \emph{include} $A_{n+1}$ (re-indexing $k\to k-1$). Together with the standalone $|A_{n+1}|$ term (the $k=1$, subset $\{n{+}1\}$ case), every subset of $\{1,\dots,n+1\}$ of every size $k=1,\dots,n+1$ is covered exactly once, with sign $(-1)^{k+1}$:
\[
\left|\bigcup_{i=1}^{n+1}A_i\right| = \sum_{k=1}^{n+1}(-1)^{k+1}\!\!\sum_{i_1<\cdots<i_k\le n+1}\!\! |A_{i_1}\cap\cdots\cap A_{i_k}|.
\]
\end{proof}

\begin{proof}[Alternative proof: per-element count]
It suffices to show that the right-hand side counts each element $a$ exactly as the left-hand side does: once if $a$ lies in the union, and $0$ times otherwise.

\textbf{Case 1} ($a\notin\bigcup_i A_i$): Then $a$ lies in no $A_i$, hence in no intersection $A_{i_1}\cap\cdots\cap A_{i_k}$. It is counted $0$ times.

\textbf{Case 2} ($a\in\bigcup_i A_i$): Let
\[
M=\{\,i : a\in A_i\,\},\qquad m=|M|\ge1.
\]
Then $a\in A_{i_1}\cap\cdots\cap A_{i_k}$ iff $\{i_1,\dots,i_k\}\subseteq M$. For each $k\le m$, exactly $\binom{m}{k}$ of the $k$-subsets of indices lie inside $M$ (none for $k>m$), and each contributes $1$ with sign $(-1)^{k+1}$. So $a$ is counted
\[
\sum_{k=1}^{m}(-1)^{k+1}\binom{m}{k}
\]
times.

\textbf{Evaluating the sum}: By Corollary~\ref{thm:altsum} (valid since $m\ge1$),
\[
0=\sum_{k=0}^{m}(-1)^k\binom{m}{k} = 1-\sum_{k=1}^{m}(-1)^{k+1}\binom{m}{k},
\]
so $a$ is counted exactly once.
\end{proof}
\begin{example}
An element lying in all of $A,B,C$ is counted $3-3+1=1$ times by the three-set formula: the case $m=3$.
\end{example}

\begin{corollary}[Dual form -- intersection via unions]
\[
\left|\bigcap_{i=1}^n A_i\right| = \sum_{k=1}^n (-1)^{k+1} \sum_{1\le i_1<\cdots<i_k\le n} |A_{i_1}\cup\cdots\cup A_{i_k}|.
\]
\end{corollary}
\begin{note}[Proof strategy]
Either repeat the induction above with $\cup\leftrightarrow\cap$ swapped throughout, or apply the theorem above to $\overline{A_1},\dots,\overline{A_n}$ and De Morgan's laws (\S\ref{sec:1-11}).
\end{note}

\subsection{Inclusion-exclusion: derangements}
\label{sec:8-4}
% CN 8.4

\begin{definition}[Derangement]
A \emph{derangement} of $\{1,\dots,n\}$ is a permutation $\sigma$ with $\sigma(i)\neq i$ for every $i$ (no fixed points).
\end{definition}

\begin{definition}[Fixed point]
For a permutation $\sigma$ of $\{1,\dots,n\}$, $i$ is a \emph{fixed point} if $\sigma(i)=i$. A derangement is a permutation with no fixed points.
\end{definition}

\begin{example}[Permutations of $\{1,2,3\}$]
Writing $\sigma$ as $(\sigma(1),\sigma(2),\sigma(3))$:
\begin{center}
\begin{tabular}{c|c}
$\sigma$ & \# fixed points \\ \hline
$(1,2,3)$ & $3$ \\
$(1,3,2)$ & $1$ \\
$(2,1,3)$ & $1$ \\
$(2,3,1)$ & $0$ \\
$(3,1,2)$ & $0$ \\
$(3,2,1)$ & $1$
\end{tabular}
\end{center}
So $D_3=2$, the derangements being $(2,3,1)$ and $(3,1,2)$.
\end{example}

\begin{theorem}[Derangement count]
\label{thm:derangecount}
\[
D_n = n!\sum_{k=0}^n \frac{(-1)^k}{k!}.
\]
\end{theorem}
\begin{proof}
Let $A_i$ = permutations of $\{1,\dots,n\}$ with $\sigma(i)=i$. Then $D_n = n! - |A_1\cup\cdots\cup A_n|$.

For any $k$ indices $i_1<\cdots<i_k$: $|A_{i_1}\cap\cdots\cap A_{i_k}|$ = permutations fixing those $k$ positions, i.e.\ permutations of the remaining $n-k$ elements: $(n-k)!$. There are $\binom{n}{k}$ such $k$-subsets.

By inclusion-exclusion:
\[
|A_1\cup\cdots\cup A_n| = \sum_{k=1}^n (-1)^{k+1}\binom{n}{k}(n-k)!.
\]
Since $\binom{n}{k}(n-k)! = \dfrac{n!}{k!}$:
\[
D_n = n! - \sum_{k=1}^n (-1)^{k+1}\frac{n!}{k!} = n! + \sum_{k=1}^n (-1)^k\frac{n!}{k!} = n!\sum_{k=0}^n \frac{(-1)^k}{k!}
\]
(the $k=0$ term contributes $n!$ back in).
\end{proof}

\begin{proposition}[Recurrence]
\label{thm:derangerec}
For $n\ge2$,
\[
D_n = (n-1)\big(D_{n-1}+D_{n-2}\big), \qquad D_0=1,\ D_1=0.
\]
\end{proposition}
\begin{proof}
\textbf{Choosing $\sigma(n)$}: In a derangement, $\sigma(n)=j$ for some $j\ne n$, giving $n-1$ choices of $j$. It remains to show that, for each fixed $j$, exactly $D_{n-1}+D_{n-2}$ derangements have $\sigma(n)=j$.

Fix $j$. Positions $1,\dots,n-1$ must be filled with the values $\{1,\dots,n\}\setminus\{j\}$, which include $n$. Split by whether position $j$ receives $n$.

\textbf{Case 1} ($\sigma(j)=n$): Then $j$ and $n$ are swapped. The other $n-2$ elements $\{1,\dots,n\}\setminus\{j,n\}$ must fill their own positions with no fixed point: $D_{n-2}$ ways.

\textbf{Case 2} ($\sigma(j)\ne n$): The constraints on positions $1,\dots,n-1$ are
\[
\sigma(i)\ne i\quad (i\ne j), \qquad \sigma(j)\ne n,
\]
because the value $j$ is already used at position $n$, so $\sigma(j)\ne j$ holds automatically. Thus the value $n$ plays exactly the role of the value $j$: its only forbidden position is $j$. Define
\[
\Phi:\ (\sigma(1),\dots,\sigma(n-1)) \longmapsto \text{same sequence, with value } n \text{ replaced by } j.
\]
\begin{itemize}
  \item $\Phi(\sigma)$ is a derangement of $\{1,\dots,n-1\}$: positions $i\ne j$ still avoid $i$, and position $j$ now avoids $j$ because it avoided $n$.
  \item $\Phi$ is invertible: given a derangement $\tau$ of $\{1,\dots,n-1\}$, replace the value $j$ by $n$ and append $\sigma(n)=j$. Since $\tau(j)\ne j$, the value $n$ does not land in position $j$, so the result lies in Case 2.
\end{itemize}
By the bijection rule (Theorem~\ref{thm:bijectionrule}), Case 2 contains $D_{n-1}$ derangements.

\textbf{Combining}: Cases 1 and 2 are disjoint and exhaustive, so each $j$ contributes $D_{n-1}+D_{n-2}$ (sum rule, \S\ref{sec:8-1}). Over the $n-1$ choices of $j$ (product rule, \S\ref{sec:8-2}) this gives the recurrence.

\textbf{Boundary values}: $D_1=0$, since the only permutation $(1)$ fixes $1$. $D_0=1$, since the empty permutation has no fixed point. Check $n=2$: $D_2=1\cdot(D_1+D_0)=1$, matching the single derangement $(2,1)$.
\end{proof}

\begin{example}[The two cases for $n=4$, $j=3$]
The derangements of $\{1,2,3,4\}$ with $\sigma(4)=3$ are:
\begin{center}
\begin{tabular}{c|c|l}
$\sigma$ & $\sigma(3)$ & case \\ \hline
$(2,1,4,3)$ & $4$ & Case 1: $\{1,2\}$ deranged as $(2,1)$ \\
$(2,4,1,3)$ & $1$ & Case 2: $\Phi$ gives $(2,3,1)$ \\
$(4,1,2,3)$ & $2$ & Case 2: $\Phi$ gives $(3,1,2)$
\end{tabular}
\end{center}
Case 2 recovers exactly the $D_3=2$ derangements of $\{1,2,3\}$, so $j=3$ contributes $1+2=3$. Over the $3$ choices of $j$, $D_4=9$.
\end{example}

\begin{example}[Small values]
\begin{center}
\begin{tabular}{c|ccccccc}
$n$ & $0$ & $1$ & $2$ & $3$ & $4$ & $5$ & $6$ \\ \hline
$D_n$ & $1$ & $0$ & $1$ & $2$ & $9$ & $44$ & $265$
\end{tabular}
\end{center}
With memoization, the recurrence computes $D_n$ using $O(n)$ arithmetic operations. Check against Theorem~\ref{thm:derangecount} at $n=4$:
\[
24\left(1-1+\frac12-\frac16+\frac1{24}\right)=9.
\]
\end{example}

\begin{corollary}[Limiting probability]
\[
\lim_{n\to\infty} \frac{D_n}{n!} = \frac{1}{e}.
\]
\end{corollary}
\begin{proof}
\[
\frac{D_n}{n!} = \sum_{k=0}^n \frac{(-1)^k}{k!},
\]
the $n$-th partial sum of the Taylor series $e^{-1}=\sum_{k=0}^\infty \dfrac{(-1)^k}{k!}$ (this partial sum converges to $e^{-1}$ as $n\to\infty$).
\end{proof}

\begin{note}[Interpretation]
The fraction of permutations with \emph{no} fixed point converges rapidly to $1/e\approx0.368$ -- already accurate to $4$ decimal places by $n=7$. This underlies the classic ``hat-check problem'': the probability nobody gets their own hat back, for $n$ people with hats mixed uniformly at random, is $D_n/n!\to1/e$, essentially independent of $n$ once $n$ is even moderately large.
\end{note}

\subsection{Selection}
\label{sec:8-5}
% CN 8.5

\begin{note}[Four kinds of selection]
\begin{center}
\begin{tabular}{lll}
 & \textbf{order matters} & \textbf{order irrelevant} \\
\textbf{with replacement} & ordered, w/ replacement (\S\ref{sec:8-6}) & unordered, w/ replacement (\S\ref{sec:8-9}) \\
\textbf{without replacement} & ordered, w/o replacement (\S\ref{sec:8-7}) & unordered, w/o replacement (\S\ref{sec:8-8})
\end{tabular}
\end{center}
\end{note}

\subsection{Ordered selection with replacement}
\label{sec:8-6}
% CN 8.6

\begin{proposition}
The number of length-$r$ sequences over an $n$-element set $A$ is
\[
|A\times\cdots\times A| = n^r
\]
(\S\ref{sec:1-14}, \S\ref{sec:1-5} for the string special case, \S\ref{sec:2-12} for counting functions).
\end{proposition}

\subsection{Ordered selection without replacement}
\label{sec:8-7}
% CN 8.7

\begin{definition}[Falling factorial, permutations]
\[
(n)_r = n(n-1)(n-2)\cdots(n-r+1) = \frac{n!}{(n-r)!}, \qquad \text{also written } {}_nP_r \text{ or } P(n,r).
\]
\end{definition}
\begin{proof}[Derivation]
Choice $k$ (for $k=1,\dots,r$) has $n-(k-1)$ options remaining (the $k-1$ prior choices excluded). By the product rule (\S\ref{sec:8-2}):
\[
n\cdot(n-1)\cdots(n-r+1).
\]
\end{proof}

\begin{note}[Counting injections]
Already the counting-injections theorem of \S\ref{sec:2-12}.
\end{note}

\subsection{Unordered selection without replacement}
\label{sec:8-8}
% CN 8.8

\begin{proposition}
\[
\binom{n}{r} = \frac{n!}{(n-r)!\,r!} = \frac{(n)_r}{r!}.
\]
\end{proposition}
\begin{note}[Dividing out the order]
Already \S\ref{sec:1-10}; each $r$-subset corresponds to $r!$ ordered sequences (\S\ref{sec:8-7}), collapsed by dividing.
\end{note}

\bigskip
\begin{theorem}[Division rule (``Second Rule of Counting'') (extra)]
\label{thm:divisionrule}
Let $A,B$ be finite and let $f:A\to B$ be \emph{$m$-to-one}, i.e.\ every $b\in B$ has exactly $m$ preimages. Then
\[
|B| = \frac{|A|}{m}.
\]
\end{theorem}
\begin{proof}
Each $a\in A$ has exactly one image, so the preimage sets $f^{-1}(b)$, $b\in B$, are pairwise disjoint with union $A$. By the sum rule (\S\ref{sec:8-1}),
\[
|A| = \sum_{b\in B}\left|f^{-1}(b)\right| = \sum_{b\in B} m = m\,|B|.
\]
\end{proof}
\begin{note}[The hypothesis is essential]
Every $b$ must have the \emph{same} number $m$ of preimages. \S\ref{sec:8-9} shows a natural situation where this fails.
\end{note}

\begin{example}[Poker hands]
Let $A$ be the set of ordered $5$-card deals, $B$ the set of $5$-card hands, and let $f$ send a deal to its set of cards. Each hand arises from exactly $5!$ deals, so $f$ is $5!$-to-one and
\[
|B| = \frac{52\cdot51\cdot50\cdot49\cdot48}{5!} = \frac{52!}{47!\,5!} = 2\,598\,960.
\]
\end{example}

\bigskip
\begin{proposition}[Pascal's identity, combinatorial proof]
\[
\binom{n}{r} = \binom{n-1}{r-1} + \binom{n-1}{r}.
\]
\end{proposition}
\begin{proof}
Fix element $x\in A$, $|A|=n$. Every $r$-subset of $A$ either contains $x$ (choose the remaining $r-1$ from $A\setminus\{x\}$: $\binom{n-1}{r-1}$ ways) or excludes $x$ (choose all $r$ from $A\setminus\{x\}$: $\binom{n-1}{r}$ ways) -- exhaustive, mutually exclusive cases.
\end{proof}

\begin{corollary}[Alternating sum]
\label{thm:altsum}
For $n\ge1$,
\[
\sum_{k=0}^n(-1)^k\binom{n}{k}=0.
\]
\end{corollary}
\begin{proof}
Let $|A|=n\ge1$ and fix $a\in A$. The map $X\mapsto X\triangle\{a\}$ (add $a$ if absent, remove it if present) is a bijection from the even-size subsets of $A$ to the odd-size ones. So there are equally many of each, i.e.\ $\sum_{k \text{ even}}\binom nk=\sum_{k \text{ odd}}\binom nk$, which is the claim.
\end{proof}
\begin{note}[Why $n\ge1$]
The hypothesis $n\ge1$ is needed: for $n=0$ the sum is $\binom00=1$. This corollary drives the combinatorial proof of inclusion--exclusion in \S\ref{sec:8-3}.
\end{note}

\begin{note}[Combinatorial proofs]
To prove an identity combinatorially, show that both sides count the same set. Either count one set in two ways (\emph{double counting}), or exhibit a bijection between the sets counted by the two sides (Theorem~\ref{thm:bijectionrule}). Pascal's identity above is a double-counting proof, and so is the next one.
\end{note}

\begin{proposition}[Hockey-stick identity (extra)]
\label{thm:hockeystick}
For $0\le k<n$,
\[
\binom{n}{k+1}=\binom{n-1}{k}+\binom{n-2}{k}+\cdots+\binom{k}{k}=\sum_{i=1}^{n-k}\binom{n-i}{k}.
\]
\end{proposition}
\begin{proof}
\textbf{Left side}: $\binom{n}{k+1}$ counts the $(k+1)$-subsets $X\subseteq\{1,\dots,n\}$.

\textbf{Right side} (classify $X$ by its minimum): Suppose $\min X=i$. Then $X=\{i\}\cup Y$, where $Y$ is a $k$-subset of $\{i+1,\dots,n\}$, a set with $n-i$ elements. So there are $\binom{n-i}{k}$ such $X$.
\begin{itemize}
  \item $Y$ needs $k$ elements, so $n-i\ge k$, i.e.\ $1\le i\le n-k$.
  \item Every $X$ has exactly one minimum, so these cases are disjoint and exhaustive.
\end{itemize}

By the sum rule (\S\ref{sec:8-1}),
\[
\binom{n}{k+1}=\sum_{i=1}^{n-k}\binom{n-i}{k}.
\]
\end{proof}
\begin{example}
$n=5$, $k=1$: $\binom52=10$, and $\binom41+\binom31+\binom21+\binom11=4+3+2+1=10$. The $2$-subsets of $\{1,\dots,5\}$ with minimum $1,2,3,4$ number $4,3,2,1$ respectively.
\end{example}
\begin{note}[Name]
In Pascal's triangle the summands lie along a diagonal, and the left-hand side sits just off the end of that diagonal, giving a hockey-stick shape.
\end{note}

\subsection{Unordered selection with replacement}
\label{sec:8-9}
% CN 8.9

\begin{definition}[Multiset]
A \emph{multiset} is a collection in which elements may repeat and order is irrelevant, e.g.\ $\{1,1,3,5\}$. A size-$r$ multiset drawn from $\{1,\dots,n\}$ is determined by its multiplicities $x_1,\dots,x_n\ge0$, where $x_1+\cdots+x_n=r$.
\end{definition}
\begin{note}[Multisets and stars and bars]
Hence the number of size-$r$ multisets from an $n$-set equals the stars-and-bars count below. This is exactly ``unordered selection with replacement''.
\end{note}

\begin{example}[Why the division rule fails here]
Choose $r=5$ pieces of fruit from $n=3$ kinds (apple $=1$, banana $=2$, orange $=3$). There are $3^5=243$ ordered selections (\S\ref{sec:8-6}). Grouping them by the resulting multiset:
\begin{center}
\begin{tabular}{l|c}
multiset & \# ordered selections \\ \hline
$\{2,2,2,2,2\}$ & $1$ \\
$\{1,2,2,2,2\}$ & $5$ \\
$\{1,1,2,2,3\}$ & $\dfrac{5!}{2!\,2!\,1!}=30$
\end{tabular}
\end{center}
The number of preimages is not constant, so Theorem~\ref{thm:divisionrule} does not apply. Indeed the answer turns out to be $21$, and $243/21$ is not an integer.
\end{example}

\begin{theorem}[Stars and bars]
The number of $(x_1,\dots,x_n)$ with each $x_i\ge0$ integer and $x_1+\cdots+x_n=r$ is
\[
\binom{r+n-1}{n-1} = \binom{r+n-1}{r}.
\]
\end{theorem}
\begin{proof}
Substitute $y_i=x_i+1\ge1$, so $\sum y_i = r+n$. Represent $r+n$ as $r+n$ ones in a row; the values $y_1,\dots,y_n$ correspond to $n-1$ barrier positions splitting them into $n$ groups, each barrier placed in one of the $r+n-1$ gaps \emph{between} consecutive ones (no barrier before the first or after the last, since each $y_i\ge1$), with at most one barrier per gap (since each $y_i\ge1$, not $\ge0$).

This is exactly choosing $n-1$ of the $r+n-1$ gap-positions, without replacement, unordered (\S\ref{sec:8-8}):
\[
\binom{r+n-1}{n-1}.
\]
\end{proof}

\begin{proof}[Alternative proof: $0/1$ strings]
Encode $(x_1,\dots,x_n)$ as the string
\[
\underbrace{0\cdots0}_{x_1}\;1\;\underbrace{0\cdots0}_{x_2}\;1\;\cdots\;1\;\underbrace{0\cdots0}_{x_n}:
\]
each unit is a $0$, and the $n-1$ separators between consecutive ``bins'' are $1$s.

\textbf{Bijection}: onto the strings of length $r+n-1$ with exactly $r$ zeros.
\begin{itemize}
  \item Every encoding has $r$ zeros and $n-1$ ones.
  \item Conversely, the $n-1$ ones cut any such string into $n$ (possibly empty) blocks of zeros, and the block lengths recover $(x_1,\dots,x_n)$.
  \item These two maps are mutually inverse.
\end{itemize}

\textbf{Counting the strings}: A string is determined by which $r$ of its $r+n-1$ positions hold $0$, an unordered selection without replacement (\S\ref{sec:8-8}). By Theorem~\ref{thm:bijectionrule},
\[
\#\{(x_1,\dots,x_n)\} = \binom{r+n-1}{r}.
\]
\end{proof}

\begin{example}
$n=5$, $r=4$: the ones split $00110110$ into the blocks $00$, (empty), $0$, (empty), $0$. So $(x_1,\dots,x_5)=(2,0,1,0,1)$, i.e.\ the multiset $\{1,1,3,5\}$. The length is $8=r+n-1$.
\end{example}

\begin{example}[Fruit salad]
$n=3$, $r=5$: there are $\binom{3+5-1}{5}=\binom75=21$ possible salads.
\end{example}

\begin{note}[Four different counts]
All four answer ``choose $r$ from $n$'' differing only in order-matters and replacement-allowed -- yet the formulas look structurally unrelated at first glance. The stars-and-bars bijection above is what connects the fourth case back to ordinary binomial coefficients.
\end{note}

